\documentclass[10pt]{article}
\usepackage[margin=0.86in]{geometry}
\usepackage[T1]{fontenc}
\usepackage{amsmath,amssymb,amsthm}
\usepackage{graphicx}
\usepackage[colorlinks=true,linkcolor=blue,citecolor=blue,urlcolor=blue]{hyperref}

\newtheorem{theorem}{Theorem}
\newtheorem{proposition}{Proposition}
\newcommand{\R}{\mathbb R}
\newcommand{\cS}{\mathcal S}

\title{The first-order heat tent square function has\\
a sharp threshold at $\sigma=1$}
\author{Run \texttt{fa\_banach\_001} --- \texttt{agent\_lane\_11}}
\date{13 August 2026}

\begin{document}
\maketitle

\noindent\textbf{Status.} Candidate full counterexample; likely valid subject
to human review. It answers negatively the all-$\sigma$ tent/Sobolev
equivalence asked by Portal and Veraar~\cite{source}.

\section{The question}

For a rough time-dependent divergence-form operator, let $\Gamma(t,s)$ be its
evolution family. The source defines
\[
 \|g\|_{T^{p,2}_\sigma}^p
 =\int_{\R^d}\left(\int_0^\infty
 \frac1{|B(x,\sqrt t)|}\int_{B(x,\sqrt t)}|g(y,t)|^2\,dy\,
 \frac{dt}{t^{1+\sigma}}\right)^{p/2}dx
\]
and asks, for $\sigma\geq0$, whether
\begin{equation}\label{eq:question}
 \|(-\Delta)^{\sigma/2}f\|_{L^p(\R^d)}
 \asymp
 \|(t,x)\mapsto\sqrt t\,\nabla\Gamma(t,0)f(x)\|_{T^{p,2}_\sigma}.
\end{equation}

\begin{center}
\includegraphics[width=0.94\textwidth]{figures/source_question_crop.png}
\end{center}

As written for every $\sigma\geq0$, the answer is no in the simplest possible
case.

\begin{theorem}\label{thm:negative}
Take identity coefficients, so that $\Gamma(t,0)=e^{t\Delta}$, and take
$p=2$. For every $\sigma\geq1$ and every nonzero
$f\in\cS(\R^d)$, the left side of \eqref{eq:question} is finite and the right
side is infinite. Consequently \eqref{eq:question} fails for every
$\sigma\geq1$ even for the ordinary heat equation.
\end{theorem}

\section{The endpoint obstruction}

At $p=2$, Fubini removes the spatial averaging in the tent norm. Therefore
\begin{align}
 &\|(t,x)\mapsto\sqrt t\,\nabla e^{t\Delta}f(x)\|_{T^{2,2}_\sigma}^{,2}
 \notag\\
 &\hspace{35mm}=
 \int_0^\infty t^{-\sigma}\|\nabla e^{t\Delta}f\|_2^2\,dt.
 \label{eq:time}
\end{align}

\begin{proof}[Proof of Theorem~\ref{thm:negative}]
The heat semigroup is strongly continuous on $H^1(\R^d)$, and hence
\[
 \nabla e^{t\Delta}f\longrightarrow\nabla f
 \quad\text{in }L^2(\R^d)\qquad(t\downarrow0).
\]
A nonzero $L^2$ Schwartz function cannot have zero gradient. Thus there are
$c,\delta>0$ such that
$\|\nabla e^{t\Delta}f\|_2^2\geq c$ for $0<t<\delta$. Formula
\eqref{eq:time} is bounded below by
\[
 c\int_0^\delta t^{-\sigma}\,dt=\infty
 \qquad(\sigma\geq1).
\]
On the other hand, $f\in\cS$ implies
$(-\Delta)^{\sigma/2}f\in L^2$ for every $\sigma\geq0$.
\end{proof}

This is not a rough-coefficient pathology: identity coefficients satisfy all
ellipticity and measurability hypotheses of the source. Nor is it an endpoint
convention: for every $\sigma>1$ the divergence is stronger.

\section{The sharp positive range for the heat model}

The same calculation identifies the exact threshold and verifies all powers.

\begin{proposition}\label{prop:positive}
For $0\leq\sigma<1$ and $f\in\cS(\R^d)$,
\begin{equation}\label{eq:exact}
 \|(t,x)\mapsto\sqrt t\,\nabla e^{t\Delta}f(x)\|_{T^{2,2}_\sigma}^{,2}
 =2^{\sigma-1}\Gamma(1-\sigma)
 \|(-\Delta)^{\sigma/2}f\|_2^2.
\end{equation}
\end{proposition}

\begin{proof}
By Plancherel, Tonelli, and $\lambda=|\xi|^2$,
\begin{align*}
 \int_0^\infty t^{-\sigma}\|\nabla e^{t\Delta}f\|_2^2dt
 &=\int_{\R^d}|\widehat f(\xi)|^2|\xi|^2
   \left(\int_0^\infty t^{-\sigma}e^{-2t|\xi|^2}dt\right)d\xi\\
 &=2^{\sigma-1}\Gamma(1-\sigma)
   \int_{\R^d}|\xi|^{2\sigma}|\widehat f(\xi)|^2d\xi.
\end{align*}
The gamma integral is finite exactly when $\sigma<1$.
\end{proof}

Thus a corrected first-order formulation can only ask about
$0\leq\sigma<1$. At and above one, a square-function kernel with additional
small-time cancellation (for example, more powers of $t(-\Delta)$) is needed.
This also shows that the all-$\sigma$ ``change of square function'' heuristic
immediately preceding the source question cannot apply to the displayed
first-order square function at $\sigma\geq1$.

\section{Verification, novelty, and limitation}

The included script checks the subcritical gamma identity and the logarithmic
and power divergences of truncated time integrals. These computations are
sanity checks; equations \eqref{eq:time}--\eqref{eq:exact} are the proof.

A bounded search on 13 August 2026 used the exact formula, paper title, and
tent/Sobolev terminology. It found both the arXiv and published source, which
repeat the same question, but no later record of this correction. The novelty
verdict is therefore \emph{apparently new within the bounded search}, not an
exhaustive bibliographic claim.

The theorem fully answers the question as quantified over all
$\sigma\geq0$. It does not decide the corrected and substantially deeper
problem for $0<\sigma<1$, rough time-dependent coefficients, and general
$p$.

\begin{thebibliography}{9}
\bibitem{source}
P. Portal and M. Veraar,
\emph{Stochastic maximal regularity for rough time-dependent problems},
Stoch. Partial Differ. Equ. Anal. Comput. \textbf{7} (2019), 541--597;
arXiv:1810.01183, p.~35;
\href{https://doi.org/10.1007/s40072-019-00134-w}{doi:10.1007/s40072-019-00134-w}.
\end{thebibliography}

\end{document}
